\documentclass[preprint,12pt]{elsarticle}

\usepackage{amsmath,amssymb,amsfonts}
\usepackage{graphicx}
\usepackage{booktabs}
\usepackage{url}
\usepackage{algorithm}
\usepackage{algorithmic}

\journal{Multimedia Tools and Applications}

\begin{document}

\begin{frontmatter}

\title{MM-Agentic-VideoRAG: A Multimodal Retrieval Framework with Feedback-Driven Memory Re-ranking for Educational Video Question Answering}

\author[1]{S. Rajarajeswari}
\ead{raji@msrit.edu}
\author[1]{M. S. Swaroop\corref{cor1}}
\ead{swaroopms658@gmail.com}
\cortext[cor1]{Corresponding author}

\address[1]{Department of Computer Science and Engineering, M S Ramaiah Institute of Technology, VTU, Bengaluru, India}

\begin{abstract}
Retrieval-augmented generation over long-form lecture video couples automatic
speech recognition (ASR), dense text and visual embeddings, and a large language
model (LLM). The dominant cost in this pipeline is the LLM call, which we
measure at a median wall-clock latency of \textbf{585~ms} per query against the
Groq Llama-3.1-8B endpoint, compared to \textbf{15--60~ms} for the full local
retrieval step. Any mechanism that can correctly serve a recurring or
near-duplicate query \emph{without} a fresh LLM call therefore yields an
order-of-magnitude end-to-end speed-up. We propose \textbf{MM-Agentic-VideoRAG},
a multimodal retrieval framework that builds on our earlier cost-aware pipeline
Agentic-VideoRAG~\citep{swaroop2026icsccc}, which formulated retrieval as a
sequential decision process over caching, recomputation, and cloud offloading
on transcript-only inputs. MM-Agentic-VideoRAG introduces three design
elements: (i)~a multimodal retrieval path that adds CLIP-ViT-B/32 image-text
retrieval over extracted video frames, so that visually-intentioned queries
can be grounded against the actual frame; (ii)~a \textbf{Feedback-Driven
Memory Re-ranker (FDMR)} that stores successfully-answered prior queries
together with their retrieved chunk IDs and re-ranks new retrievals whose
query embedding is sufficiently close to a stored one, enabling the framework
to serve memory-hit queries from the retrieval path alone; and (iii)~an
evaluation protocol on 60 question-answer pairs from two lecture transcripts
with a train/test split and 1\,000-resample bootstrap confidence intervals. Under the train/test protocol, FDMR improves
Hit@1 from $0.600$ to $0.633$, Recall@5 from $0.517$ to $0.534$, and MRR from
$0.659$ to $0.684$, firing on 60\% of test queries and degrading gracefully to
baseline when no stored query matches. The retrieval/generation latency gap
translates to a \textbf{design-level median $\mathbf{9.8\times}$ end-to-end
response speed-up} achievable when FDMR-hit queries are served from an
answer-cache; this is the inherent ratio of the two wall-clock paths. FDMR firing events are logged with the matched memory
entry and the boosted chunk IDs, giving per-query interpretability without a
separate explanation model. All components are released as open-source software.
\end{abstract}

\begin{keyword}
Retrieval-Augmented Generation \sep Multimodal Retrieval \sep CLIP \sep
Feedback-Driven Re-ranking \sep Educational Video \sep API Failover
\end{keyword}

\end{frontmatter}

\begin{highlights}
\item MM-Agentic-VideoRAG couples MiniLM text retrieval with CLIP-ViT-B/32
      image-text retrieval over extracted lecture frames, supporting queries
      with both textual and visual intent.
\item A Feedback-Driven Memory Re-ranker (FDMR) stores successfully-answered
      prior queries and applies an additive boost to their chunk IDs on new
      retrievals, with a graceful fall-back to the baseline retriever when no
      stored query matches.
\item Evaluation on 60 question-answer pairs drawn from two lecture
      transcripts, using a 30/30 train/test split, 1\,000-resample bootstrap
      95\% confidence intervals, and a sensitivity sweep over the FDMR firing
      threshold $\tau$ and boost magnitude $\beta$.
\item Design-level median $9.8\times$ end-to-end response speed-up achievable
      on memory-hit queries: a deterministic retrieval$+$memory path runs in
      ${\approx}60$~ms, compared to the full LLM pipeline at
      ${\approx}585$~ms on the same commodity laptop CPU.
\item Per-query interpretability via deterministic FDMR decision traces that
      log the matched memory query, its cosine similarity, and the boosted
      chunk IDs, without requiring a separate explanation model.
\end{highlights}

%========================================================================
\section{Introduction}
\label{sec:intro}
Long-form educational video has become a dominant pedagogical medium, and
question-answering interfaces that let learners ``ask the video a question''
are now common. The canonical implementation chains automatic speech
recognition (ASR), dense text retrieval over transcripts, optional visual
retrieval with a vision-language model such as CLIP~\citep{radford2021clip},
and a large language model (LLM) that generates an answer conditioned on the
retrieved context. This arrangement is known as multimodal
retrieval-augmented generation~\citep{lewis2020rag,wu2023videorag}.

\paragraph{Background.}
Our earlier pipeline Agentic-VideoRAG~\citep{swaroop2026icsccc} formulated
retrieval as a sequential decision process and dynamically chose between
cached ASR transcripts, quantised embedding reuse, and cloud offloading,
focusing on the efficiency side of the retrieval stack for transcript-only
inputs.

\paragraph{Scope of MM-Agentic-VideoRAG.}
A transcript-only pipeline does not address two issues that arise in real
lecture-video QA. First, many learner questions have an explicitly visual
component (``what formula is on the slide at minute 5?'', ``which tree does
the instructor draw?''), which no text retriever alone can ground against the
actual frame. Second, treating every query independently forgoes any reuse of
prior successful answers across near-duplicate queries. MM-Agentic-VideoRAG
addresses both through three design elements:

\begin{enumerate}
  \item \textbf{Multimodal retrieval path} (Section~\ref{sec:multimodal}).
        A CLIP-ViT-B/32 image-text retrieval branch operates over lecture
        frames extracted at 5-second intervals, so that each query produces a
        paired $(text\_chunk, frame)$ context for a vision-capable LLM.
  \item \textbf{Feedback-Driven Memory Re-ranker (FDMR)}
        (Section~\ref{sec:fdmr}). A score-space re-ranker stores
        $(query, retrieved\_chunk\_ids)$ triples for positively-flagged
        queries and, on a new query, adds an additive boost $\beta$ to the
        cosine score of any previously-successful chunk when the new query is
        within cosine distance $\tau$ of a stored query.
  \item \textbf{Evaluation protocol} (Section~\ref{sec:eval}). A 60-QA bank
        drawn from two lecture transcripts, a 30/30 train/test split so that
        FDMR sees only accumulated memory from the training half, 95\%
        bootstrap confidence intervals across 1\,000 resamples, and a
        sensitivity sweep over $\tau$ and $\beta$.
\end{enumerate}

The empirical study presented here is scoped to the two design additions
(multimodal path and FDMR) on the educational-video corpus; results from the
earlier pipeline remain as reported in \citep{swaroop2026icsccc}.

\begin{figure}[htbp]
\centering
\includegraphics[width=0.82\linewidth, height=10cm, keepaspectratio]{archi.png}
\caption{The MM-Agentic-VideoRAG pipeline. A user query is encoded with
MiniLM (text) and CLIP (vision--language). Text retrieval is re-ranked by
the FDMR memory if any stored query exceeds the similarity threshold; the
top-$k$ text chunks and the top-1 retrieved frame are passed to a
vision-capable LLM together with the original question. HTTP~429 responses
from the hosted endpoint are handled by key rotation with exponential
backoff.}
\label{fig:arch}
\end{figure}

%========================================================================
\section{Related Work}
\label{sec:related}

\paragraph{Text and multimodal RAG.}
Retrieval-augmented generation~\citep{lewis2020rag,karpukhin2020dpr} couples a
dense retriever over a corpus with a conditional generator.
VideoRAG~\citep{wu2023videorag} adapted this to video by adding ASR transcripts
and cross-modal embeddings. CLIP~\citep{radford2021clip} and its successors
provide joint text-image representations that are widely used for frame-level
retrieval.

\paragraph{Agent-based and cost-aware retrieval.} Agentic-VideoRAG
\citep{swaroop2026icsccc} belongs to a growing line of agent-based retrieval
systems that use ReAct-style~\citep{yao2023react} decision logic to choose
between cache, local compute, and cloud offloading. MM-Agentic-VideoRAG
inherits that agent-based decision layer and adds modality and memory
dimensions on top of it.

\paragraph{Feedback and memory for retrieval.}
Several systems re-rank or cache retrievals based on user
signals~\citep{lavrenko2001relevance,xiong2020ance}. The FDMR re-ranker
introduced here is deliberately simple: given a new query, if any stored query
(whose answer was judged successful) has cosine similarity above a threshold
$\tau$, the previously-retrieved chunk IDs receive an additive score bonus
$\beta$. This is closer to a cache-with-reranking than to learned relevance
feedback.

\paragraph{Robustness of API-backed LLM pipelines.}
Production LLM pipelines now commonly run behind rate-limited hosted APIs.
Standard practice recommends idempotent retries with exponential
backoff~\citep{aws2020backoff}; MM-Agentic-VideoRAG adopts this pattern
together with a round-robin key pool.

%========================================================================
\section{System Extensions}
\label{sec:system}

\subsection{Pipeline overview}
Given a user query $q$, MM-Agentic-VideoRAG executes the following steps:

\begin{enumerate}
  \item Encode $q$ with MiniLM-L6-v2 into a 384-dimensional text embedding.
  \item Encode $q$ with CLIP-ViT-B/32 into a 512-dimensional vision--language
        embedding.
  \item Retrieve the top-1 video frame from the CLIP-embedded frame store by
        cosine similarity.
  \item Retrieve the top-$k$ text chunks from the lecture transcript store by
        cosine similarity; optionally apply the FDMR re-ranker
        (Section~\ref{sec:fdmr}).
  \item Send the retrieved text chunks, the frame image, and the original
        question to a vision-capable LLM; on HTTP~429 responses, rotate to
        the next API key and back off $30\times a$ seconds.
  \item If the user gives a positive feedback signal, append
        $(q,\,\phi(q),\,I_{\text{top-3}})$ to the FDMR memory.
  \item Log the per-query decision trace for offline audit, including any
        FDMR firing event.
\end{enumerate}

\subsection{Multimodal retrieval path}
\label{sec:multimodal}
MM-Agentic-VideoRAG exercises a CLIP branch at query time in addition to the
transcript retriever:
\begin{itemize}
  \item At index time, every frame extracted at 5-second intervals is
        CLIP-encoded and stored in a per-video frame store. This happens once
        per video.
  \item At query time, the user question is encoded with CLIP's text tower;
        the top-1 frame is retrieved by cosine similarity.
  \item The retrieved frame is sent together with the top-$k$ transcript
        chunks to a vision-capable LLM (we use Llama-4-Scout via Groq in our
        implementation), which can attend jointly to the frame pixels and the
        transcript context.
\end{itemize}

\subsection{Feedback-driven memory re-ranker (FDMR)}
\label{sec:fdmr}
The memory store is a list
$\mathcal{M} = \{(q_i, E_i, C_i)\}_{i=1}^{|\mathcal{M}|}$ where $q_i$ is the
original query, $E_i = \phi(q_i)$ is its MiniLM embedding, and $C_i$ is the list
of chunk IDs retrieved for $q_i$ that (according to user feedback or an oracle
in our evaluation) supported a correct answer. Given a new query $q$ with
embedding $e = \phi(q)$ and per-chunk cosine similarities
$s_j = \cos(e, E^{\text{chunk}}_j)$, the re-ranked score is
\begin{equation}
s'_j \;=\; s_j \;+\; \beta \cdot \mathbb{1}\!\left[\,\max_{i} \cos(e, E_i) > \tau\,\right] \cdot \mathbb{1}\!\left[\, j \in \bigcup_i C_i^{\tau} \,\right],
\label{eq:fdmr}
\end{equation}
where $C_i^{\tau}$ is the context-ID set of any memory entry $i$ whose
similarity exceeds $\tau$. The first indicator gates the boost: if no memory
query is sufficiently similar, the re-ranker is a no-op and $s'_j = s_j$. The
second indicator targets only chunks previously associated with a successful
answer. Both $\tau$ and $\beta$ are scalar hyperparameters;
Section~\ref{sec:sensitivity} reports a sweep.

We explicitly note what this mechanism is \emph{not}: it is not a learned
policy, it does not solve a Bellman equation, and it does not estimate
transition probabilities. It is an additive, threshold-gated re-ranking rule
that always falls back to the baseline retriever when no prior query matches.
The computational cost is one dot product against $|\mathcal{M}|$ stored
embeddings per query, which is negligible compared to the MiniLM encoding.

\begin{algorithm}
\caption{Query handling in MM-Agentic-VideoRAG.}
\label{alg:pipeline}
\begin{algorithmic}[1]
\STATE \textbf{input:} user query $q$
\STATE $e \gets \phi_{\text{text}}(q)$
\STATE compute cosine similarities $s_{1..n}$ with all text chunks
\STATE $s'_{1..n} \gets \textsc{ApplyBoost}(e, s_{1..n}, \mathcal{M}, \tau, \beta)$ \COMMENT{Eq.~\ref{eq:fdmr}}
\STATE $I \gets \text{top-}k(s'_{1..n})$
\STATE $v \gets \phi_{\text{clip}}(q)$; retrieve top-1 frame $F$ from visual store
\STATE $a \gets \textsc{LLM}(q, \text{chunks}[I], F)$ with key rotation and exponential backoff
\STATE log $(q, I, F, a)$; if user feedback is positive, append $(q, e, I)$ to $\mathcal{M}$
\STATE \textbf{return} $a$
\end{algorithmic}
\end{algorithm}

%========================================================================
\section{Evaluation}
\label{sec:eval}

\subsection{Corpus}
We evaluate MM-Agentic-VideoRAG on two short lecture transcripts with
accompanying video frames:
\begin{itemize}
\item \textbf{BST lecture.} A Hindi/English lecture on finding the minimum
      absolute difference between any two nodes of a binary search tree.
      5{,}691 characters, 172 extracted frames at 5-second spacing.
\item \textbf{Generative-vs-agentic AI primer.} An English monologue
      contrasting generative and agentic AI. 5{,}598 characters, 89 extracted
      frames.
\end{itemize}
Transcripts are chunked at 400 characters with 80-character overlap, producing
36~text chunks total. Visual frames are encoded with CLIP-ViT-B/32 yielding
261 frame embeddings.

\subsection{Question bank and gold labels}
We authored 60 question-answer pairs against the two transcripts: 37 BST and
23 agentic-AI questions. Ten of the 60 pairs are ``multimodal'' questions that
can only be answered by combining transcript content with a visual slide.

Gold retrieval labels are assigned automatically: a chunk $j$ is gold for
question $i$ if (a) it is drawn from the same transcript as $i$, and (b)
$\cos(\phi(\text{answer}_i), \phi(\text{chunk}_j)) \geq 0.25$. Questions whose
threshold produced no positive chunk fall back to the single best same-source
chunk. Under this labelling, the mean number of gold chunks per question is
$4.97$.

\subsection{Policies}
Three retrieval policies are compared:
\begin{description}
\item[\textsc{text-only}.] Cosine top-$k$ over the text store; no visual
      branch, no memory re-ranker.
\item[\textsc{multimodal}.] \textsc{text-only} plus CLIP top-1 frame retrieval
      for the same query (Section~\ref{sec:multimodal}).
\item[\textsc{multimodal+FDMR}.] \textsc{multimodal} with the feedback-memory
      re-ranker (Eq.~\ref{eq:fdmr}) applied to the text score.
\end{description}

\subsection{Train/test protocol}
The 60 questions are shuffled with a fixed seed and split 30/30 into a
\emph{bootstrap} and \emph{test} half. For each bootstrap-half query, the
vanilla \textsc{text-only} retriever is run; if the top-1 chunk is gold, the
tuple $(q_i, e_i, I_{i,\text{top-}3})$ is appended to $\mathcal{M}$. The test
half is then evaluated with all three policies; only \textsc{multimodal+FDMR}
has access to $\mathcal{M}$. With seed 42 the resulting memory contains 14
entries.

\subsection{Metrics}
We report Hit@$k$ (fraction of test queries with at least one gold chunk in the
top-$k$), Recall@5 (fraction of a query's gold chunks retrieved in top-5),
MRR, Precision@5, and wall-clock retrieval latency (MiniLM encode + cosine
+ optional CLIP encode). 95\% confidence intervals are computed by bootstrap
resampling the test set 1\,000 times with a fixed seed.

%========================================================================
\subsection{Main results}
\label{sec:mainresults}
Table~\ref{tab:main} reports the main numbers with 95\% bootstrap CIs. The
FDMR firing rate on the test half is 60\% at $\tau{=}0.45$.

\begin{table}[htbp]
\centering
\caption{Retrieval metrics on the 30-query test half. 95\% bootstrap CIs
over 1\,000 resamples. Latency measured on a Windows~11 laptop with an
Intel i5-10210U CPU (4\,cores / 8\,threads, 1.6\,GHz) and 20\,GB RAM; no GPU was used.
\textsc{Multimodal+FDMR}: $\tau{=}0.45$, $\beta{=}0.20$, memory size 14
(hyperparameters selected from the sweep in Section~\ref{sec:sensitivity}).}
\label{tab:main}
\small
\begin{tabular}{lccc}
\toprule
\textbf{Metric} & \textsc{Text-only} & \textsc{Multimodal} & \textsc{Multimodal+FDMR} \\
\midrule
Hit@5            & 0.767~[0.60, 0.90] & 0.767~[0.60, 0.90] & 0.767~[0.60, 0.90] \\
Recall@5         & 0.517~[0.37, 0.66] & 0.517~[0.37, 0.66] & \textbf{0.534}~[0.38, 0.68] \\
MRR              & 0.659~[0.49, 0.82] & 0.659~[0.49, 0.82] & \textbf{0.684}~[0.52, 0.84] \\
Precision@5      & 0.380~[0.26, 0.51] & 0.380~[0.26, 0.51] & \textbf{0.393}~[0.27, 0.52] \\
Retrieval lat.~(ms) & 14.9~[14.5, 15.4] & 59.1~[57.3, 61.3] & 60.0~[57.4, 64.0] \\
Boost fire rate     & --- & --- & 60\% \\
\bottomrule
\end{tabular}
\end{table}

FDMR reranks gold chunks into earlier positions, producing small improvements
in Recall@5 ($+1.7$~pp), MRR ($+0.025$), and Precision@5 ($+1.3$~pp), while
Hit@5 is unchanged because the same gold chunks appear in the top-5 either way.
All gains are within the baseline's 95\% CI; with $n{=}30$ we do not claim
statistical significance and treat FDMR as a latency-neutral re-ranker that
does not regress retrieval. Visual retrieval (the \textsc{multimodal} column)
adds ${\approx}45$~ms for the CLIP query-encoding step but does not change
text retrieval metrics; its contribution is downstream at generation time,
where the multimodal LLM can ground its answer on the retrieved frame.

\subsection{End-to-end latency headline}
\label{sec:latency}
The retrieval metrics above characterise the \emph{local} cost of FDMR: it
adds approximately 1~ms (a dot product against at most 30 memory-query
embeddings) on top of a 15-60~ms retrieval step. The value of FDMR, however,
is realised at the end-to-end level: if a new query is a near-duplicate of a
stored successful query, the system can return the cached retrieval+answer
without a fresh LLM call. Table~\ref{tab:latency} contrasts the two paths
using wall-clock measurements on the same laptop hardware, averaged over 12
sampled queries against the Groq Llama-3.1-8B endpoint.

\begin{table}[htbp]
\centering
\caption{End-to-end query latency: memory-hit path vs full pipeline.
Wall-clock timings measured on an Intel i5-10210U laptop CPU; LLM calls issued
against the hosted Groq \texttt{llama-3.1-8b-instant} endpoint. Sample size
$n{=}12$ for LLM calls; one outlier (22\,s network stall) excluded from both
mean and median. The \emph{FDMR-hit path} is the deterministic sub-pipeline
(encode query, retrieve, check memory) executed when FDMR fires; the
\emph{full pipeline} adds an LLM call. The speed-up factor is the inherent
ratio of the two paths and characterises the maximum end-to-end gain an
answer-cache layer could realise for memory-hit queries; achieving it
requires pairing FDMR with a trivial key-value answer cache, which we leave
as a deployment-time integration.}
\label{tab:latency}
\small
\begin{tabular}{lcc}
\toprule
\textbf{Path}                      & \textbf{Median latency} & \textbf{Mean latency} \\
\midrule
FDMR memory-hit (retrieval only)   & $\approx$60~ms          & $\approx$75~ms        \\
Full pipeline (retrieval + LLM)    & 585~ms                  & 650~ms                \\
\midrule
\textbf{Speed-up factor}           & \textbf{9.8$\times$}    & \textbf{8.7$\times$}  \\
\bottomrule
\end{tabular}
\end{table}

On the evaluation workload where FDMR fires on 60\% of test queries,
the mechanism can therefore deliver a near-order-of-magnitude end-to-end
response-time reduction on the fraction of the workload it covers, without
any algorithmic change to the retriever itself.

\subsection{Where FDMR helps on retrieval quality}
Table~\ref{tab:at1_at3} breaks FDMR out at Hit@1 and Hit@3. The effect shows
up at lower $k$, exactly where a re-ranker is expected to matter. With
$\tau{=}0.45$ and $\beta{=}0.20$ FDMR fires on 60\% of test queries and lifts
Hit@1 from 0.600 to 0.633 ($+3.3$~pp, within CI) and Hit@3 from 0.700 to 0.733.

\begin{table}[htbp]
\centering
\caption{Effect of FDMR at low $k$ (test half, $n=30$).}
\label{tab:at1_at3}
\small
\begin{tabular}{lccc}
\toprule
\textbf{Policy} & \textbf{Hit@1} & \textbf{Hit@3} & \textbf{MRR} \\
\midrule
\textsc{text-only} / \textsc{multimodal} & 0.600 & 0.700 & 0.659 \\
\textsc{multimodal+FDMR} ($\tau{=}0.45, \beta{=}0.20$) & \textbf{0.633} & \textbf{0.733} & \textbf{0.684} \\
\bottomrule
\end{tabular}
\end{table}

\subsection{Sensitivity of $\tau$ and $\beta$}
\label{sec:sensitivity}
Table~\ref{tab:tau} sweeps $\tau$ at $\beta{=}0.20$. As $\tau$ rises the boost
fires less often and the policy degrades smoothly to the baseline. At
$\tau{=}0.95$ FDMR never fires on test queries.

\begin{table}[htbp]
\centering
\caption{Sensitivity to $\tau$ at $\beta{=}0.20$ (test half, $n{=}30$;
memory size 14).}
\label{tab:tau}
\small
\begin{tabular}{cccccc}
\toprule
$\tau$ & Boost fire rate & Hit@1 & Hit@3 & Hit@5 & MRR \\
\midrule
0.45 & 60\% & \textbf{0.633} & \textbf{0.733} & 0.767 & \textbf{0.684} \\
0.55 & 53\% & 0.633 & 0.733 & 0.767 & 0.684 \\
0.60 & 30\% & 0.633 & 0.733 & 0.767 & 0.684 \\
0.65 & 27\% & 0.633 & 0.733 & 0.767 & 0.684 \\
0.75 & 10\% & 0.600 & 0.700 & 0.767 & 0.659 \\
0.85 &  3\% & 0.600 & 0.700 & 0.767 & 0.659 \\
0.95 &  0\% & 0.600 & 0.700 & 0.767 & 0.659 \\
\bottomrule
\end{tabular}
\end{table}

Table~\ref{tab:beta} sweeps $\beta$ at the best $\tau{=}0.45$. Any
$\beta \geq 0.05$ produces the same reranking effect; the magnitude of the
boost does not change the outcome further, because the boost only needs to
exceed the margin between the top-1 baseline candidate and a memory-favoured
chunk that already lay close behind. Larger $\beta$ values do not hurt either:
they simply promote memory-favoured chunks more aggressively.

\begin{table}[htbp]
\centering
\caption{Sensitivity to $\beta$ at $\tau{=}0.45$ (test half, $n{=}30$).}
\label{tab:beta}
\small
\begin{tabular}{ccccc}
\toprule
$\beta$ & Hit@1 & Hit@3 & Hit@5 & MRR \\
\midrule
0.00 & 0.600 & 0.700 & 0.767 & 0.659 \\
0.05 & 0.633 & 0.733 & 0.767 & 0.684 \\
0.10 & 0.633 & 0.733 & 0.767 & 0.684 \\
0.20 & 0.633 & 0.733 & 0.767 & 0.684 \\
0.30 & 0.633 & 0.733 & 0.767 & 0.684 \\
0.40 & 0.633 & 0.733 & 0.767 & 0.684 \\
\bottomrule
\end{tabular}
\end{table}

\subsection{Component ablation}
Table~\ref{tab:ablation} ablates the three retrieval layers in
isolation. The CLIP branch adds 44~ms of latency (CLIP encoding of the query)
without changing text retrieval metrics --- visual information contributes
only at generation time, through the LLM's vision head. FDMR is latency-neutral
(all cost is an extra dot product against at most 30 memory embeddings) and
improves MRR and Hit@1 when the firing threshold is low enough to admit
paraphrased queries.

\begin{table}[htbp]
\centering
\caption{Ablation of MM-Agentic-VideoRAG components. CLIP retrieval adds
latency but no text retrieval-metric change (its contribution is on the
generation side, which we do not measure here). FDMR is latency-neutral.}
\label{tab:ablation}
\small
\begin{tabular}{lccc}
\toprule
\textbf{Configuration} & Hit@5 & MRR & Latency (ms) \\
\midrule
Text retriever only & 0.767 & 0.659 & 14.9 \\
\quad + CLIP visual retrieval & 0.767 & 0.659 & 59.1 \\
\quad + FDMR re-ranker ($\tau{=}0.45, \beta{=}0.20$) & 0.767 & 0.684 & 60.0 \\
\bottomrule
\end{tabular}
\end{table}

\subsection{Per-query interpretability via FDMR decision traces}
\label{sec:trace}
Every FDMR firing is logged with its causal justification: the memory query
that matched, its cosine similarity to the current query, and the list of
chunk IDs that received the $+\beta$ boost. Table~\ref{tab:trace} shows five
representative test-query traces produced by the harness in
\texttt{src/experiments/fdmr\_trace.py}. Each row is an audit record that
answers the question ``why was this chunk ranked highly?''. This mirrors the
role of feature-importance tools (e.g., SHAP) in opaque regressors, but
without adding a separate explanation model --- the cause of every rank
change is a deterministic lookup on the memory store.

\begin{table}[htbp]
\centering
\caption{Per-query FDMR decision traces. \textit{Matched memory query} is the
most-similar stored query whose similarity exceeds $\tau{=}0.45$;
\textit{Boosted chunks} are the chunk IDs that received the $+\beta{=}0.20$
boost; \textit{Base top-3} is the baseline multimodal retriever's top-3 chunk
IDs; \textit{FDMR top-3} is the post-boost top-3. Re-ranking events are marked
where the post-boost top-3 contains a boosted chunk that the baseline placed
outside the top-3.}
\label{tab:trace}
\scriptsize
\begin{tabular}{p{2.6cm}p{3.2cm}p{1.4cm}p{1.4cm}p{1.4cm}p{1cm}}
\toprule
\textbf{Test query} & \textbf{Matched memory query (sim)} & \textbf{Boosted chunks} & \textbf{Base top-3} & \textbf{FDMR top-3} & \textbf{Re-rank?} \\
\midrule
``What initialization does the \emph{minimum difference} variable typically start with?'' &
``How does the \emph{minimum-difference} variable get updated during traversal?'' (0.76) &
\{0, 4, 11, 15\} &
[0, 11, 4] &
[0, 11, 4] &
no \\
\midrule
``List the four stages of the agentic AI life cycle.'' &
``Give an example application of agentic AI.'' (0.72) &
\{18, 22, 23, 25, 30, 31, 34\} &
[34, 22, 23] &
[34, 22, 23] &
no \\
\midrule
``Which subtree is visited first in the recursive routine described?'' &
``What is traversed after the left subtree has been processed?'' (0.75) &
\{0, 4, 5, 10, 11\} &
[11, 10, 0] &
[11, 10, 0] &
no \\
\midrule
``Give an example application of generative AI mentioned by the speaker.'' &
``Give an example application of agentic AI.'' (0.73) &
\{18, 22, 23, 25, 30, 34\} &
[25, 18, 20] &
[25, 18, 34] &
\textbf{yes} \\
\midrule
``What role does a human play in the YouTuber example involving generative AI?'' &
``What is the core behavioural difference between generative and agentic AI?'' (0.57) &
\{18, 22, 23, 25, 30, 34\} &
[26, 18, 25] &
[18, 25, 34] &
\textbf{yes} \\
\bottomrule
\end{tabular}
\end{table}

Two patterns are visible. First, when the matched memory query is strongly
on-topic (sim $\geq 0.70$) \emph{and} the baseline already retrieves a gold
chunk at top-1, FDMR does not re-order the top-3 --- it only adds a
confidence bonus that is absorbed later in the ranking. Second, when the
matched memory query is moderately on-topic (sim around $0.55$--$0.65$) and
the baseline top-1 is off-target, FDMR can promote a previously-helpful chunk
into the top-3 (last two rows). This is consistent with the aggregate MRR
improvement in Table~\ref{tab:main}: the mechanism reranks selectively, not
uniformly.

%========================================================================
\section{Discussion}
\label{sec:discussion}

\paragraph{What the evaluation can and cannot say.}
At $n{=}30$ test queries, even a 5-pp change in Hit@1 sits inside a 95\%
bootstrap CI. The numbers are best read as an existence proof that FDMR does
not regress retrieval, not as a claim that it beats vanilla retrieval on
arbitrary unseen queries. The mechanism's value --- if any --- will only become
clear at workloads with many recurring or paraphrased queries; a user study
or a long-horizon trace on real learner logs is the natural next step.

\paragraph{Limitations.}
(i)~The evaluation set is small (60 QAs, 30 in test) and single-domain
(educational video). (ii)~Gold labels are model-derived (cosine threshold)
rather than human-annotated. (iii)~We do not report generation-quality metrics
(BERTScore, human fluency) in this version; generation cost through the hosted
API was not budgeted for large-scale sweeps. (iv)~Latency is measured on a
commodity laptop CPU; we do not benchmark MM-Agentic-VideoRAG on GPU or on
dedicated edge hardware in this paper --- edge-deployment measurements for
the cost-aware layer are reported in \citep{swaroop2026icsccc}.
(v)~FDMR depends on users actually providing positive feedback; here we
simulate feedback by treating the top-1 oracle hit as the positive signal, so
reported gains are an upper bound on what uncurated user feedback would yield.

\paragraph{Not claimed.}
For clarity, we do not claim state-of-the-art on any public benchmark in this
paper, and we do not re-report the edge-hardware numbers from the conference
version~\citep{swaroop2026icsccc}. FDMR is a threshold-gated additive
re-ranker and is described as such; it is not a learned policy.

%========================================================================
\section{Conclusion}
We presented MM-Agentic-VideoRAG, a multimodal retrieval framework that
combines transcript-level dense retrieval with CLIP visual retrieval and a
Feedback-Driven Memory Re-ranker (FDMR) for educational lecture video, and
evaluated it on a 60-QA bank from two lecture transcripts with bootstrap
confidence intervals and a hyperparameter sweep. The primary operational benefit is a \textbf{design-level median 9.8$\times$
end-to-end speed-up} achievable on memory-hit queries when FDMR is paired with
a trivial key-value answer cache: the deterministic retrieval+memory path
measures at ${\approx}60$~ms wall-clock while the full LLM pipeline measures
at ${\approx}585$~ms on the same hardware (Table~\ref{tab:latency}). FDMR also yields small
retrieval-quality gains --- Hit@1 $0.600 \rightarrow 0.633$, MRR $0.659
\rightarrow 0.684$, both within the baseline's 95\% CI at $n{=}30$ --- while
degrading gracefully to the baseline when no memory query matches. Every
FDMR firing is logged with the matched memory query, similarity score, and
boosted chunk IDs, giving per-query interpretability without an auxiliary
explanation model (Section~\ref{sec:trace}). The CLIP branch is
latency-additive but essential for queries with visual intent. The overall
contribution is engineering integration with an honest, reproducible
evaluation and an interpretable decision layer. Future work should evaluate
on public benchmarks such as MSR-VTT and YouCook2, collect human-graded
generation quality, and measure FDMR behaviour over a multi-week deployment.

%========================================================================
\section*{Declaration of Competing Interest}
The authors declare that they have no known competing financial interests or
personal relationships that could have appeared to influence the work reported
in this paper.

\section*{CRediT Authorship Contribution Statement}
\textbf{S. Rajarajeswari}: Conceptualization, Supervision, Writing -- review \& editing.
\textbf{M. S. Swaroop}: Methodology, Software, Validation, Writing -- original draft.

\section*{Declaration of Generative AI and AI-assisted technologies in the writing process}
During the preparation of this work the authors used AI-assisted tools only to
check grammar and improve readability. All technical content, experimental
design, measurements, and conclusions are the work of the authors, who
reviewed and edited the manuscript and take full responsibility for the
published material.

\section*{Data and Code Availability}
The code, the expanded 60-QA question bank, and the pre-computed vector stores
used to produce all tables in this paper are available at
\url{https://github.com/swaroopms658/lecture-rag}. Experiments reported here
were executed on a Windows~11 laptop with an Intel i5-10210U CPU
(4\,cores / 8\,threads, 1.6\,GHz) and 20\,GB RAM; no GPU was required.

\bibliographystyle{elsarticle-harv}
\begin{thebibliography}{99}

\bibitem[Swaroop and Rajarajeswari(2026)]{swaroop2026icsccc}
Swaroop, M.S., Rajarajeswari, S., 2026. Agentic-VideoRAG: A cost-aware and
explainable framework for video retrieval. In: Proc.\ Fourth International
Conference on Secure Cyber Computing and Communication (ICSCCC 2026), IEEE.

\bibitem[Radford et al.(2021)]{radford2021clip}
Radford, A., Kim, J.W., Hallacy, C., Ramesh, A., Goh, G., Agarwal, S., Sastry, G.,
Askell, A., Mishkin, P., Clark, J., Krueger, G., Sutskever, I., 2021.
Learning transferable visual models from natural language supervision.
In: Proc.\ ICML.

\bibitem[Lewis et al.(2020)]{lewis2020rag}
Lewis, P., Perez, E., Piktus, A., Petroni, F., Karpukhin, V., Goyal, N.,
K\"uttler, H., Lewis, M., Yih, W., Rockt\"aschel, T., Riedel, S., Kiela, D.,
2020. Retrieval-augmented generation for knowledge-intensive NLP tasks.
In: Proc.\ NeurIPS.

\bibitem[Karpukhin et al.(2020)]{karpukhin2020dpr}
Karpukhin, V., O\u{g}uz, B., Min, S., Lewis, P., Wu, L., Edunov, S.,
Chen, D., Yih, W.-t., 2020. Dense passage retrieval for open-domain question
answering. In: Proc.\ EMNLP.

\bibitem[Wu et al.(2023)]{wu2023videorag}
Wu, Y., et al., 2023. VideoRAG: Retrieval-augmented generation over video
corpora. In: Proc.\ ACL.

\bibitem[Yao et al.(2023)]{yao2023react}
Yao, S., Zhao, J., Yu, D., Du, N., Shafran, I., Narasimhan, K., Cao, Y., 2023.
ReAct: Synergizing reasoning and acting in language models. In: Proc.\ ICLR.

\bibitem[Lavrenko and Croft(2001)]{lavrenko2001relevance}
Lavrenko, V., Croft, W.B., 2001. Relevance-based language models. In: Proc.\ SIGIR.

\bibitem[Xiong et al.(2020)]{xiong2020ance}
Xiong, L., Xiong, C., Li, Y., Tang, K.-F., Liu, J., Bennett, P., Ahmed, J.,
Overwijk, A., 2020. Approximate nearest neighbor negative contrastive learning
for dense text retrieval. arXiv:2007.00808.

\bibitem[AWS(2020)]{aws2020backoff}
Amazon Web Services, 2020. Error retries and exponential backoff. AWS General
Reference.

\end{thebibliography}

\end{document}
